\chapter*{Abstract}
\addcontentsline{toc}{chapter}{Abstract}

Phishing websites use misleading URLs to imitate trusted services and obtain
sensitive information from users. Detection models trained only on historical
URLs may perform well on familiar data but lose accuracy when URL structures,
sources, or attack patterns change. This project investigates whether phishing
URLs generated by Generative Adversarial Networks (GANs) can improve the
generalisation of URL-based classifiers. Traditional and deep-learning
classifiers were trained as real-only baselines, including Random Forest,
Support Vector Machine, Convolutional Neural Network, Recurrent Neural Network,
and Long Short-Term Memory models. Several text-generation approaches were
then studied, and generated URLs were checked for validity, uniqueness,
diversity, overlap with real data, and their ability to challenge real-only
classifiers. Controlled augmentation experiments added fixed percentages of
synthetic phishing URLs while keeping the real training rows and evaluation
sets unchanged. A matched comparison also trained SeqGAN and MS-HardGAN-C on
the exact same 84,280 phishing rows; MS-HardGAN-C retained the better downstream
F1 and false-negative balance. The final system combines character-level CNN learning with 35
lexical URL features, multi-source real training, score calibration, and a
three-model ensemble. On a final test of 6,000 unseen URLs containing 5\%
phishing, the selected 10\% GAN-augmented model achieved 98.23\% accuracy,
84.15\% precision, 79.67\% recall, and 81.85\% F1-score. Its false-positive
rate was 0.79\% and its false-negative rate was 20.33\%. The model was also
exported to a browser extension that performs local URL analysis and presents
warning levels without uploading visited URLs. A later browser-policy audit
showed that blanket trust for shared hosting could hide phishing URLs. Replacing
it with a non-blocking shared-content warning reduced FNR from 21.33\% to
13.33\% on the audited snapshot without changing its measured FPR of 0.69\%.
Overall, carefully filtered GAN augmentation improved detection under some test
conditions, but generated data cannot replace diverse real data or remove all
effects of distribution shift.

\clearpage
